\subsection{Witt's Theorem.}

Given two vector spaces E, E' over A equipped respectively with two sesquilinear forms \(\Phi\) , \(\Phi'\) (resp. with two quadratic forms Q, Q'), one calls a metric homomorphism from E to E' any linear map \(u\) from E to E' such that \(\Phi'(u(x), u(y)) = \Phi(x, y)\) (resp. Q'(u(x)) = Q(x)) for \(x \in E\) , \(y \in E\) . If E and E' have the same finite dimension and if \(\Phi\) (resp. Q) is nondegenerate, every metric homomorphism \(u\) from E to E' is an isomorphism, since \(u(x) = 0\) implies \(\Phi(x, y) = 0\) for all \(y \in E\) , hence \(x = 0\) ; thus \(u\) is injective, hence bijective since E and E' have the same finite dimension.

THEOREM 1 (Witt). — Let E and E' be two finite-dimensional vector spaces, equipped respectively with two nondegenerate ε-hermitian forms Φ and Φ' satisfying condition (T) of no. 2 (resp. with two nondegenerate quadratic forms Q and Q'), and isomorphic for these structures. Given any subspace F of E, every injective metric homomorphism of F into E' extends to a metric isomorphism of E onto E'.

Using the given isomorphism from E onto \(\mathbf{E}'\) , one sees that it suffices to show that every injective metric homomorphism \(u\) from F into E extends to a metric automorphism of E. Note that if, for \(i = 1, 2\) , \(\mathbf{F}_i\) is a subspace of E and \(u_i\) a metric homomorphism from \(\mathbf{F}_i\) into E, such that \(\mathbf{F}_1 \cap \mathbf{F}_2 = \{0\}\) and \(\Phi(u_1(x_1), u_2(x_2)) = \Phi(x_1, x_2)\) for \(x_i \in \mathbf{F}_i\) ( \(i = 1, 2\) ), then the homomorphism \(\nu : x_1 + x_2 \to u_1(x_1) + u_2(x_2)\) of \(\mathbf{F}_1 + \mathbf{F}_2\) into E which extends \(u_1\) and \(u_2\) is metric: indeed, for any \(x_i, y_i\) in \(\mathbf{F}_i\) ( \(i = 1, 2\) ), the expansion of each of the expressions \(\Phi(x_1 + x_2, y_1 + y_2)\) and \(\Phi(u_1(x_1) + u_2(x_2), u_1(y_1) + u_2(y_2))\) (resp. \(\mathrm{Q}(x_1 + x_2)\) and \(\mathrm{Q}(u_1(x_1) + u_2(x_2))\) ) contains four (resp. three) terms equal each to each according to the hypotheses made. Moreover, if \(u_1\) and \(u_2\) are injective and if \(u_1(\mathbf{F}_1) \cap u_2(\mathbf{F}_2) = \{0\}\) , then \(\nu\) is injective.

\begin{enumerate}
\def\labelenumi{\arabic{enumi})}
\tightlist
\item
  Let us first prove Witt's theorem in the case where the set of fixed points of u is a hyperplane U of F. The set of vectors of the form \(u(x) - x\) with \(x \in F\) is then a line D. If \(F'\) is a subspace orthogonal to D such that \(F' \cap F = F' \cap u(F) = \{0\}\) , one will have \(\Phi(u(x), y) = \Phi(x, y)\) for \(x \in F\) and \(y \in F'\) ; our initial remark therefore applies to u and to the identity map of \(F'\) into E, showing that u extends to \(F + F'\) while leaving the points of \(F'\) fixed; the set of vectors of the form \(u(x) - x (x \in F + F')\) is still the line D. Now we have, for \(x \in F, y \in F\)
\end{enumerate}

\[
\Phi (u (x), u (y) - y) = \Phi (u (x), u (y)) - \Phi (u (x), y) = \Phi (x - u (x), y),
\]

which, when \(x \in \mathrm{U}\) (that is, when \(u(x) = x\) ), shows that \(x \in \mathrm{D}^0\) ; in other words, we have \(\mathrm{U} \subset \mathrm{D}^0\) . We distinguish two cases:

\begin{enumerate}
\def\labelenumi{\alph{enumi})}
\item
  \(F \subset D^{0}\) . Formula (5) shows that \(u(F)\) is not contained in \(D^{0}\) , hence \(F \cap D^{0} = u(F) \cap D^{0} = U\) . One can then take for \(F'\) a supplement of U in \(D^{0}\) ; since \(F + F'\) contains the hyperplane \(D^{0}\) and is distinct from it, we have \(F + F' = E\) , and we have found in this case the desired extension of u to E.
\item
  \(\mathbf{F} \subset \mathbf{D}^0\) . Formula (5) shows that \(u(\mathbf{F}) \subset \mathbf{D}^0\) , and hence that
\end{enumerate}

D ⊂ D⁰ ; the line D is therefore isotropic (resp. singular, since we have Q(u(x) - x) = Q(u(x)) - Φ(x, u(x)) + Q(x) = 2Q(x) - Φ(x, x) = 0 for x ∈ F). We shall show that, under these conditions, there exists a subspace F' of D⁰ which is a supplement of F and of u(F) in D⁰. This is immediate if F = u(F). Otherwise, let x and y be vectors such that x ∈ F, x ∉ U, y ∈ u(F), y ∉ U; we then have F = U + Ax, u(F) = U + Ay, and F does not contain x + y, otherwise y = (x + y) - x would belong to F ∩ u(F) = U; we see similarly that x + y does not belong to u(F); thus the line A(x + y) is a supplement of F and of u(F) in the subspace F + u(F); it then suffices to set F' = A(x + y) + G, where G is a supplement of F + u(F) in D⁰. This being so, we have F + F' = u(F) + F' = D⁰, and, in this case, what was said at the beginning of 1) shows that there exists an extension of u to the hyperplane D⁰ of E, and that D⁰ is stable for this extension.

We are therefore reduced to the case where F is the hyperplane D⁰ and where u is an automorphism of F. Let us prove that, for every z ∈ E, there exists z' ∈ E such that

\[
\Phi (u (x), z ^ {\prime}) = \Phi (x, z)\tag{6}
\]

for all \(x \in \mathbf{F}\) ; indeed the linear form \(x \to \Phi(u^{-1}(x), z)\) on \(\mathbf{F}\) is the restriction of a linear form on \(\mathbf{E}\) , a form which is of the type \(x \to \Phi(x, z')\) since \(\Phi\) is nondegenerate; hence (6) is valid. Moreover, if \(z \notin \mathbf{F}\) there exists a vector \(z' \in \mathbf{E}\) satisfying (6) and such that \(\Phi(z', z') = \Phi(z, z)\) (resp. \(Q(z') = Q(z)\) ). Indeed, formula (6) remains valid if one adds to \(z'\) an element \(u(y) - y\) ( \(y \in \mathbf{F}\) ) of \(\mathbf{D}\) since \(\mathbf{F} = \mathbf{D}^0\) , and Lemma 1 of No. 2 allows us to conclude since \(z\) is not orthogonal to \(\mathbf{D}\) Our initial remark then shows that there exists a metric homomorphism \(v\) of \(\mathbf{F} + A z = E\) into \(E\) which extends \(u\) and which transforms \(z\) to \(z'\) . Since \(\Phi\) is nondegenerate, \(v\) is the metric automorphism of \(E\) sought.

\begin{enumerate}
\def\labelenumi{\arabic{enumi})}
\setcounter{enumi}{1}
\tightlist
\item
  In the general case, we will argue by induction on \(r = \dim F\) The case \(r = 0\) is trivial. Let then \(r > 0\) , that is to say \(F \neq \{0\}\) and let U be a hyperplane of F. The restriction \(u_0\) of \(u\) to U extends, by the induction hypothesis, to a metric automorphism \(v_0\) of E. If \(v_0\) extends \(u\) , the theorem is proved. Otherwise, U is the set of invariant elements by \(v_{0}^{-1}u\) , and there exists, by 1), a metric automorphism \(v_{1}\) of E extending \(v_{0}^{-1}u\) . The automorphism \(v_{0}v_{1}\) is then the desired extension of u. QED.
\end{enumerate}

COROLLARY 1. --- Let, for \(i=1,2\) , \(\mathbf{E}_{i}\) be a finite-dimensional vector space, \(\Phi_{i}\) a nondegenerate \(\varepsilon\) -hermitian form on \(\mathbf{E}_{i}\) satisfying (T) (resp. \(\mathbf{Q}_{i}\) a nondegenerate quadratic form on \(\mathbf{E}_{i}\) ), \(\mathbf{E}_{i}^{\prime}\) and \(\mathbf{E}_{i}^{\prime\prime}\) two orthogonal subspaces of \(\mathbf{E}_{i}\) of which \(\mathbf{E}_{i}\) is the direct sum. If the forms \(\Phi_{1}\) and \(\Phi_{2}\) (resp. \(\mathbf{Q}_{1}\) and \(\mathbf{Q}_{2}\) ) are equivalent, and if their restrictions to \(\mathbf{E}_{1}^{\prime}\) and \(\mathbf{E}_{2}^{\prime}\) are equivalent, then so are their restrictions to \(\mathbf{E}_{1}^{\prime\prime}\) and \(\mathbf{E}_{2}^{\prime\prime}\) .

Indeed, let \(u\) be a metric isomorphism of \(\mathbf{E}_{1}^{\prime}\) onto \(\mathbf{E}_{2}^{\prime}\) . By th. 1, \(u\) extends to a metric isomorphism \(\nu\) of \(\mathbf{E}_{1}\) onto \(\mathbf{E}_{2}\) . Since \(\Phi_{i}\) is nondegenerate, \(\mathbf{E}_{i}^{\prime\prime}\) is the orthogonal of \(\mathbf{E}_{i}^{\prime}\) in \(\mathbf{E}_{i}\) , hence \(\nu\) maps \(\mathbf{E}_{1}^{\prime\prime}\) on \(\mathbf{E}_{2}^{\prime\prime}\) . QED.

COROLLARY 2. --- Under the hypotheses of Theorem 1, the group of metric automorphisms of E acts transitively on the totally isotropic (resp. totally singular) subspaces of E of a given dimension. Moreover, if F is a totally isotropic (resp. totally singular) subspace of E, then every bijective linear map from F onto F is induced by a metric automorphism of E.

COROLLARY 3. --- Let Q be a nondegenerate quadratic form on a finite-dimensional vector space E over an algebraically closed field A. The group of metric automorphisms of E acts transitively on the non-isotropic subspaces of E of a given dimension.

This follows immediately from th. 1 and cor. 2 of prop. 3.

Exercises. -- 1) a) Let K be a field of characteristic 2, J : \(\xi \to \bar{\xi}\) an involutive antiautomorphism of K, Z the center of K. Show that if the restriction of J to Z is not the identity, every element \(\mu\) of K such that \(\bar{\mu} = \mu\) is of the form \(\lambda + \bar{\lambda}\) (note that there is an element \(\rho \neq 0\) in Z which can be written in the form \(\zeta + \bar{\zeta}\) , with \(\zeta \in \mathbb{Z}\) ); every Hermitian form on a vector space over K then satisfies condition (T).

\begin{enumerate}
\def\labelenumi{\alph{enumi})}
\setcounter{enumi}{1}
\tightlist
\item
  Give examples of fields of characteristic 2 admitting an involutive antiautomorphism \(\xi \to \overline{\xi}\) distinct from the identity map, and for which there exist elements \(\mu = \bar{\mu}\) which are not of the form \(\lambda + \bar{\lambda}\) (cf. Chap. VIII, § 11, Exerc. 4).
\end{enumerate}

\begin{enumerate}
\def\labelenumi{\arabic{enumi})}
\setcounter{enumi}{1}
\tightlist
\item
  Let A be a field, E a vector space over A, \(\Phi\) (resp. Q) a nondegenerate \(\varepsilon\) -hermitian form on E, satisfying condition (T) (resp. a nondegenerate quadratic form on E, \(\Phi\) then denoting the symmetric bilinear form associated with Q).
\end{enumerate}

\begin{enumerate}
\def\labelenumi{\alph{enumi})}
\item
  Show that for a plane P ⊂ E to be isotropic (resp. singular) and not totally isotropic (resp. not totally singular), it is necessary and sufficient that it contain only one isotropic (resp. singular) line (cf. exerc. 14 e)).
\item
  We assume that dim \(E \geqslant 3\) and that there exist isotropic vectors in E \(\neq 0\) Show that if P is a non-totally isotropic plane in E, there exists a non-isotropic subspace \(V \subset E\) of dimension 3, containing isotropic vectors \(\neq 0\) , and such that \(P \subset V\) .
\end{enumerate}

\begin{enumerate}
\def\labelenumi{\arabic{enumi})}
\setcounter{enumi}{2}
\item
  With the same hypotheses as in exerc. 2, show that if dim E ≥ 3, every isotropic line in E is the intersection of two non-isotropic planes.
\item
  The hypotheses are those of exerc. 2, and it is further assumed that E is of finite dimension.
\end{enumerate}

\begin{enumerate}
\def\labelenumi{\alph{enumi})}
\item
  If the index v of \(\Phi\) (resp. Q) is \(\geqslant 1\) , show that for every isotropic (resp. singular) vector \(a \neq 0\) in E, there exists a basis \((e_i)\) of E consisting of isotropic (resp. singular) vectors, such that \(e_1 = a\) (cf. exerc. 14 e)).
\item
  Let V, W be two totally isotropic (resp. totally singular) subspaces of the same dimension \(r \leqslant v\) ; show that there exist two maximal totally isotropic (resp. totally singular) subspaces \(V_{1}\) , \(W_{1}\) , such that \(V \subset V_{1}\) , \(W \subset W_{1}\) , and \(V_{1} \cap W_{1} = V \cap W\) . (If U = V ∩ W, reason in \(U^{0}/U\) ).
\item
  Let V, W, \(V_{1}\) , \(W_{1}\) be four totally isotropic (resp. totally singular) subspaces of the same dimension, such that \(V + W\) and \(V_{1} + W_{1}\) are non-isotropic. Show that there exists a metric automorphism u of E such that \(u(V) = V_{1}\) and \(u(W) = W_{1}\) .
\item
  Let \(f\) be a linear form on E, \(\alpha\) an element of A of the form \(\lambda + \varepsilon\bar{\lambda}\) (resp. an element of A). We consider the sesquilinear form on E
\end{enumerate}

\[
(x, y) \rightarrow \Phi_ {1} (x, y) = \Phi (x, y) + f (x) \alpha \overline {{f (y)}}
\]

(resp. the quadratic form

\[
x \rightarrow \mathrm{Q} _ {1} (x) = \mathrm{Q} (x) + \alpha (f (x)) ^ {2}).
\]

Show that if \(\Phi_{1}\) (resp. \(Q_{1}\) ) is nondegenerate and if \(\nu_{1}\) denotes its index, we have \(|\nu_{1} - \nu| \leqslant 1\) .

\begin{enumerate}
\def\labelenumi{\arabic{enumi})}
\setcounter{enumi}{4}
\tightlist
\item
  \begin{enumerate}
  \def\labelenumii{\alph{enumii})}
  \tightlist
  \item
    Let B be a ring, \(\xi \to \overline{\xi}\) an involutive antiautomorphism of B, \(\varepsilon\) an element of the center of B such that \(\varepsilon \varepsilon = 1\) . Show that if \(\beta\) is an invertible element of B such that \(\beta + \varepsilon \overline{\beta} \neq 0\) , there exists an invertible element \(\mu \neq 1\) in B, such that \(\mu (\beta + \varepsilon \overline{\beta}) \mu = \beta + \varepsilon \overline{\beta}\) (Show that one can take \(\mu\) such that \(\mu \beta \overline{\mu} = \beta\) ).
  \end{enumerate}
\end{enumerate}

\begin{enumerate}
\def\labelenumi{\alph{enumi})}
\setcounter{enumi}{1}
\tightlist
\item
  Let A be a field, and let E be a vector space over A, \(\Phi\) a nondegenerate ε-hermitian sesquilinear form on E, satisfying (T). Show that if Φ is not alternating, then for every non-isotropic hyperplane H of E, there exists a metric automorphism of E, distinct from the identity, leaving every element of H invariant (use a)).
\end{enumerate}

\begin{enumerate}
\def\labelenumi{\arabic{enumi})}
\setcounter{enumi}{5}
\item
  Given the hypotheses of exercise 2, let \(a\) be an isotropic vector \(\neq 0\) in E (resp. a nonsingular isotropic vector (note that such vectors exist only if A has characteristic 2)). Let \(\lambda \in \mathbf{A}\) such that \(\lambda + \varepsilon \bar{\lambda} = 0\) (resp. \(\lambda = (\mathrm{Q}(a))^{-1}\) ); show that the transvection \(x \to x + \Phi(x, a) \lambda a\) (chap. II, § 6, exerc. 7) is a metric automorphism of E; conversely.
\item
  Under the hypotheses of Exercise 2, let G be the group of metric automorphisms of E. Show that the only semilinear bijections of E onto itself that commute with every element of G are the homotheties of E, except in the following three cases: dim E = 2, G is the group of metric automorphisms corresponding to a quadratic form of index 1 on E, and A is one of the three fields \(\mathbf{F}_{2}, \mathbf{F}_{3}\) or \(\mathbf{F}_{4}\) (Use Exercises 5, 6, and 3; examine separately the case of a quadratic form on a vector space of dimension 2).
\end{enumerate}

*8) Let A be a field, E a vector space of finite dimension \textgreater0 over A, Φ a nondegenerate ε-hermitian sesquilinear form on E, satisfying (T). Let M(Φ) be the group of multipliers of the similarities of E for Φ (§ 6, no. 5).

\begin{enumerate}
\def\labelenumi{\alph{enumi})}
\item
  Let \(V_{1}\) , \(V_{2}\) two vector subspaces of E, of the same dimension, and let \(\Phi_{1}\) , \(\Phi_{2}\) be the restrictions of \(\Phi\) to \(V_{1}\) , \(V_{2}\) respectively. For there to exist a similarity u such that \(u(V_{1}) = V_{2}\) , it is necessary and sufficient that there exist \(\alpha \in M(\Phi)\) such that \(\Phi_{2}\) is equivalent to \(\alpha\Phi_{1}\) (use Witt's theorem).
\item
  Let (F, F', G) be a Witt decomposition of E (no. 2), and let \(\Phi_0\) be the restriction of \(\Phi\) to the non-isotropic subspace G. Show that one has M( \(\Phi\) ) = M( \(\Phi_0\) ) if G ≠ \{0\}. (Use Witt's theorem and prop. 2 of no. 2).
\item
  Show that if the index v of Φ is such that dim E = 2v, M(Φ) is the group of elements ζ ≠ 0 of the center of A such that ζ̄ = ζ. If dim E = 2v + 1, M(Φ) is the group of elements of the form ρρ̄, where ρ runs through the multiplicative group of elements ≠ 0 of the center of A (use Witt's theorem). (Cf. § 10, exerc. 18).*
\end{enumerate}

*9) Let A be a commutative field, E a vector space over A, Q a nondegenerate quadratic form on E. One calls a similarity for Q any automorphism \(u\) of E such that there exists an element \(\alpha \neq 0\) of A for which \(Q(u(x)) = \alpha Q(x)\) for every \(x \in E\) ; \(u\) is then also a similarity for the bilinear form associated with Q. Assuming the dimension of E finite and \(>0\) , state and prove for the similarities relative to Q the analogues of the results of exerc. 8.*

*10) Let A be a field, E a vector space over A of dimension \(\geqslant 2\) , \(\Phi_1\) (resp. \(\Phi_2\) ) a nondegenerate sesquilinear form \(\varepsilon_1\) -hermitian (resp. \(\varepsilon_2\) -hermitian) on E for an involutive antiautomorphism \(J_1\) (resp. \(J_2\) ) of A, satisfying condition (T). Show that if the group of metric automorphisms of E for \(\Phi_1\) is a subgroup of the group of similarities for \(\Phi_2\) , there exists \(\alpha \in A\) such that \(\Phi_2 = \Phi_1\alpha\) (use Exercises 5 b) and 6).

Prove the analogous property when A is assumed commutative, and that \(\Phi_{1}\) and \(\Phi_{2}\) are replaced in the statement by two nondegenerate quadratic forms \(Q_{1}, Q_{2}\) on E.*

\begin{enumerate}
\def\labelenumi{\arabic{enumi})}
\setcounter{enumi}{10}
\tightlist
\item
  Let A be a ring, J an involutive antiautomorphism of A, and E an A-module admitting a finite basis \((e_{i})\) , \(\Phi\) a nondegenerate \(\varepsilon\) -hermitian form on E, R the matrix of \(\Phi\) with respect to \((e_{i})\) ; the group of metric automorphisms of \(\Phi\) is identified with the group G of invertible matrices U such that \(^{t}U.R.U^{j}=R\) .
\end{enumerate}

\begin{enumerate}
\def\labelenumi{\alph{enumi})}
\tightlist
\item
  Assume that there exists a matrix \(P\) such that \(R = {}^t P + \varepsilon P^J\) . Show that for every matrix \(S\) such that \({}^t S + \varepsilon S^J = 0\) , and that \(P + S\) is invertible, \(U = (^{t}P^{J} - \varepsilon^{J}S)^{-1}(P + S)\) belongs to G, and \(\varepsilon I + U\) is invertible. Conversely (show that for every matrix \(U \in G\) such that \(\varepsilon I + U\) is invertible, we have
\end{enumerate}

\[
\varepsilon (\varepsilon I + ^ {t} U) ^ {- 1} R + \varepsilon^ {\mathrm{J}} R (\varepsilon^ {\mathrm{J}} I + U ^ {\mathrm{J}}) ^ {- 1} = R).
\]

\begin{enumerate}
\def\labelenumi{\alph{enumi})}
\setcounter{enumi}{1}
\tightlist
\item
  Show that the condition of a) is satisfied when \(\Phi\) satisfies condition (T). Case where in A the equation \(2\xi = \alpha\) admits one and only one solution for every \(\alpha \in \mathbf{A}\) .
\end{enumerate}

¶ 12) Let A be a commutative field, E a vector space of finite dimension n over A, Φ a nondegenerate ε-hermitian sesquilinear form on E.

\begin{enumerate}
\def\labelenumi{\alph{enumi})}
\item
  Let \(u\) be an endomorphism of E; show that if the \(r_i\) ( \(1 \leqslant i \leqslant m\) ) are the similarity invariants of \(u\) (chap. VII, § 5, no. 1, def. 1), the similarity invariants of the adjoint \(u^*\) of \(u\) with respect to \(\Phi\) are the polynomials \(\bar{r}_i\) ( \(1 \leqslant i \leqslant m\) ), where \(\bar{r}_i\) is deduced from \(r_i\) by applying to each coefficient the automorphism J (cf. chap. VII, § 5, exerc. 2). For every monic irreducible polynomial \(p \in \mathrm{A}[\mathrm{X}]\) dividing the minimal polynomial of \(u\) , let \(\mathrm{F}_k(u, p)\) be the kernel of \((p(u))^k\) in E, and let \(\mathrm{F}(u, p)\) be the union of the \(\mathrm{F}_k(u, p)\) for all integers \(k > 0\) . Show that if \(p\) and \(q\) are two distinct monic irreducible polynomials dividing the minimal polynomial of \(u\) , the subspaces \(\mathrm{F}(u, p)\) and \(\mathrm{F}(u^*, \bar{q})\) are orthogonal (use Bezout's identity). Finally, if G is a vector subspace of E such that \(u(\mathrm{G}) \subset \mathrm{G}\) , one has \(u^*(\mathrm{G}^0) \subset \mathrm{G}^0\) .
\item
  We assume that \(uu^{*} = u^{*}u\) (the case where one says that \(u\) is a normal endomorphism for \(\Phi\) ; cf.§ 7, no. 3); show that one then has \(u^{*}(F_{k}(u,p)) \subset F_{k}(u,p)\) for all \(k\) and consequently \(u^{*}(F(u,p)) \subset F(u,p)\) . If one sets \(G(p,\bar{q}) = F(u,p) \cap F(u^{*},\bar{q})\) , show that E is the direct sum of the subspaces \(G(p,\bar{q})\) , and that \(G(p,\bar{q})\) and \(G(p_{1},\bar{q}_{1})\) are orthogonal if \(p \neq q_{1}\) or if \(p_{1} \neq q\) ; in particular \(G(p,\bar{q})\) is totally isotropic if \(p \neq q\) . Show that \(G(p,\bar{p})\) is reduced to 0 or non-isotropic, and that if \(p \neq q\) , no nonzero vector of \(G(p,\bar{q})\) is orthogonal to \(G(q,\bar{p})\) (using the fact that \(\Phi\) is nondegenerate); deduce that if \(p \neq q\) , \(G(p,\bar{q})\) and \(G(q,\bar{p})\) are totally isotropic subspaces of the same dimension, and that \(G(p,\bar{q}) + G(q,\bar{p})\) is non-isotropic.
\item
  Assume that J is not the identity or that A does not have characteristic 2, and that \(u^{*} = u\) . Let M be the set of non-isotropic subspaces \(\mathbf{M} \subset \mathrm{G}(p, \overline{p})\) stable under \(u\) (hence submodules of the A{[}X{]}- module \(E_{u}\) (chap. VII, § 5, no. 1)). Show that if M is a minimal element of M, then M is an indecomposable submodule of \(E_{u}\) (chap. VII, § 4, no. 7). (Assume that M is a direct sum of an indecomposable submodule \(M_{1}\) and of a submodule \(M_{2} \neq \{0\}\) , the minimal polynomials \(p^{h}\) and \(p^{k}\) of the restrictions of u to \(M_{1}\) and \(M_{2}\) respectively being such that \(h \geqslant k\) . Note then that \(M_{1}\) is necessarily isotropic and that every \(z \neq 0\) in \(M_{1}\) such that \(p(u).z = 0\) is orthogonal to \(M_{1}\) (using the fact that every submodule of \(M_{1}\) is monogenic); writing that \(z = (p(u))^{h-1}.x\) and that z is not orthogonal to \(M_{2}\) , and deduce that one necessarily has k = h. Then show that there exists an indecomposable submodule \(N_{2}\) of \(M_{2}\) such that \(p^{h}\) be the minimal polynomial of the restriction of u to \(N_{2}\) , and that \(M_{1} + N_{2}\) be non-isotropic; conclude from this that \(M_{2} = N_{2}\) . Finally, if \(y \in M_{2}\) is not orthogonal to z, consider the submodule P of M generated by \(w = x + \lambda y\) , where \(\lambda \in A\) , and show that one can take \(\lambda\) such that P is non-isotropic, by proving that one has \(\Phi((p(u))^{h-1}.w,w) \neq 0\) ; which leads to a contradiction).
\end{enumerate}

\(d)\) Deduce from \(c)\) that \(\mathrm{G}(p,\bar{p})\) is a direct sum of indecomposable submodules \(\mathrm{H}_{i}\) pairwise orthogonal. If \(p^{h}\) is the minimal polynomial of the restriction of \(u\) to \(\mathrm{H}_{i}\) , and if \(d\) is the degree of \(p\) , show that there exists in \(\mathrm{H}_{i}\) a totally isotropic subspace of dimension \(d\) . {[}h/2{]}. Case where E contains no isotropic vector \(\neq 0\) (cf.§ 7, no. 3).

\begin{enumerate}
\def\labelenumi{\alph{enumi})}
\setcounter{enumi}{4}
\item
  State and prove the properties analogous to those of c) and d) when one has \(u^{*} = -u\) or \(u^{*}u = 1\) .
\item
  Give an example where n = 4, Φ is symmetric and of index 2, p = \(\bar{p} = X - 1\) , u is normal, E = G(p, \(\bar{p}\) ), but E is not a direct sum of minimal submodules of M, and where there exists an eigenvector of u that is not an eigenvector of \(u^{*}\) (cf.§ 7, no. 3).
\end{enumerate}

\begin{enumerate}
\def\labelenumi{\arabic{enumi})}
\setcounter{enumi}{12}
\item
  The hypotheses are those of Exercise 2, and one further assumes that E admits a countable basis \((e_n)\) . Let F be a totally isotropic (resp. totally singular) subspace of E such that \(F^{00} = F\) . Show that there exists a totally isotropic (resp. totally singular) subspace \(F'\) such that: \(1^\circ F \cap F' = \{0\}\) ; \(2^\circ\) there exists a basis \((a_m)_{m \in I}\) of F and a basis \((a'_m)_{m \in I}\) of \(F'\) (I interval of N starting at 0) such that \(\Phi(a_i, a_j') = \delta_{ij}\) for every pair of indices; \(3^\circ (F + F')^{00} = F + F'\) and E is the direct sum of \(F + F'\) and of \(G = (F + F')^0\) . (Construct by induction an increasing sequence \((L_n)\) of non-isotropic subspaces, with union E, such that dim \(L_{n+1} - \dim L_n = 2\) , and apply prop. 2 of no. 2 to each of the \(L_n\) ; to construct this sequence, one will consider, for every \(n\) , the smallest integer \(k\) such that \(e_k \notin L_n\) , and one will use exerc. 9 b) of § 1).
\item
  Let A be a field of characteristic 2, E a vector space of finite dimension n over A, Φ a nondegenerate Hermitian form on E, not necessarily satisfying condition (T).
\end{enumerate}

\begin{enumerate}
\def\labelenumi{\alph{enumi})}
\item
  Show that the set V of \(x \in E\) such that \(\Phi(x, x)\) is of the form \(\alpha + \bar{\alpha}\) is a vector subspace of E.
\item
  Let \(V_{1}=V\cap V^{0}\) , \(q=\dim V_{1}\) , \(V_{2}\) a supplement of \(V_{1}\) with respect to V, \(V_{3}\) a supplement of \(V_{1}\) with respect to \(V^{0}\) . Show that there exists a basis \((e_{i})_{1\leqslant i\leqslant2q}\) of \((\mathrm{V}_{2}+\mathrm{V}_{3})^{0}=\mathrm{V}_{2}^{0}\cap\mathrm{V}_{3}^{0}\) such that the vectors \(e_1, \ldots, e_q\) form a basis of \(V_1\) and that one has \(\Phi(e_i, e_{q+j}) = \delta_{ij}\) for \(1 \leqslant i \leqslant q\) , \(1 \leqslant j \leqslant q\) .
\item
  Let \(\mathbf{G}(\Phi)\) be the group of metric automorphisms of E (for \(\Phi\) ). Show that for every \(u \in \mathbf{G}(\Phi)\) , one has \(u(x) = x\) for all \(x \in V^0\) .
\item
  For every \(u \in \mathbf{G}(\Phi)\) , one has \(u(V) = V\) ; let \(u_{v}\) be the restriction of \(u\) to \(V\) , and let \(G_{v}\) be the group formed by the \(u_{v}\) . Show that: 1° the kernel of the homomorphism \(u \to u_{v}\) of \(\mathbf{G}(\Phi)\) onto \(G_{v}\) is commutative; 2° if \(\Phi_{2}\) is the restriction of \(\Phi\) to \(V_{2}\) and \(\mathbf{G}(\Phi_{2})\) the group of metric automorphisms of \(V_{2}\) for \(\Phi_{2}\) , there exists a homomorphism from \(G_{v}\) onto \(\mathbf{G}(\Phi_{2})\) whose kernel is commutative (use \(b\) ) and \(c\) ).
\item
  We assume that A is commutative and J is the identity; let E be a vector space of dimension 3 over A, \((e_{i})_{1\leqslant i\leqslant3}\) a basis of E, \(\Phi\) the nondegenerate symmetric form on E whose matrix with respect to \((e_{i})\) is
\end{enumerate}

\[
\left( \begin{array}{c c c} 1 & 0 & 0 \\ 0 & 0 & 1 \\ 0 & 1 & 0 \end{array} \right).
\]

Show that all isotropic vectors in E are contained in the hyperplane spanned by \(e_{2}\) and \(e_{3}\) (cf. exercise 4(a)). Give an example of a non-isotropic plane containing only one isotropic line (cf. exercise 2(a)). Show that there exists no automorphism \(u \in \mathbf{G}(\Phi)\) such that \(u(e_{1}) = e_{1} + e_{2}\) , even though one has \(\Phi(e_{1}, e_{1}) = \Phi(e_{1} + e_{2}, e_{1} + e_{2})\) .

\section{Special properties of alternating bilinear forms}

\subsection{Reduction of alternating bilinear forms.}

THEOREM 1. --- Let A be a principal (commutative) ring, E a free A-module of finite dimension n, and \(\Phi\) an alternating bilinear form on E. Then there exists a basis \((e_i)_{1 \leqslant i \leqslant n}\) of E and an even integer \(2r \leqslant n\) , such that

\[
1 ^ {0} \quad \Phi (e _ {1}, e _ {2}) = \alpha_ {1}, \Phi (e _ {3}, e _ {4}) = \alpha_ {2}, \dots , \Phi (e _ {2 r - 1}, e _ {2 r}) = \alpha_ {r}
\]

where the \(\alpha_{i}\) are elements \(\neq 0\) of A, and where \(\alpha_{i}\) divides \(\alpha_{i+1}\) for \(i=1,\ldots,r-1\) .

2° All the other elements \(\Phi(e_{i}, e_{j})\) where \(i \leqslant j\) are zero.

The ideals \(\mathrm{A}\alpha_{i}\) ( \(i=1,\ldots,r\) ) are uniquely determined by the preceding conditions. The submodule \(\mathrm{E}^{0}\) of E orthogonal to E is generated by \(e_{2r+1},\ldots,e_{n}\) .

We shall proceed by induction on the dimension n of E. The theorem is obvious for n = 0. If \(\Phi = 0\) , the theorem is also obvious; we may therefore assume \(\Phi \neq 0\) . Let f denote the linear map \(d_{\Phi}\) from E into \(\mathbf{E}^*\) associated on the right to \(\Phi\) ( \(\S 1, \text{no. 1}\) ); then \(f(\mathbf{E})\) is a submodule not reduced to 0 of the module \(\mathbf{E}^*\) , which is a free module of dimension \(n\) . Let \(\mathrm{A}\alpha_{1}\) be the largest invariant factor of \(f(\mathbf{E})\) with respect to \(\mathbf{E}^*\) (chap. VII, \(\S 4, \text{no. 2}\) , th. 1); we know (loc. cit.) that there exists a basis \((e_{1}', a_{2}', \ldots, a_{n}')\) of \(\mathbf{E}^*\) and an element \(f(e_{2}) \in f(\mathbf{E})\) such that \(f(e_{2}) = \alpha_{1}e_{1}'\) . Let \((e_{1}, a_{2}, \ldots, a_{n})\) the basis of E (identified with the bidual \(\mathbf{E}^{**}\) ) dual of \((e_{1}', a_{2}', \ldots, a_{n}')\) ; we have

\[
\Phi (e _ {1}, e _ {2}) = - \Phi (e _ {2}, e _ {1}) = \left\langle e _ {1}, f (e _ {2}) \right\rangle = \alpha_ {1}.\tag{1}
\]

Let P be the submodule \(Ae_{1} + Ae_{2}\) of E. We shall see that E is the direct sum of \(Ae_{1}\) , \(Ae_{2}\) and of the submodule \(P^{0}\) orthogonal to P. For this it suffices to prove that, for every \(x \in E\) , there exist elements \(\xi_{1}\) , \(\xi_{2}\) of A, uniquely determined and such that \(x - \xi_{1}e_{1} - \xi_{2}e_{2} \in P^{0}\) that is, such that

\[
\Phi (e _ {1}, x - \xi_ {1} e _ {1} - \xi_ {2} e _ {2}) = 0, \quad \Phi (e _ {2}, x - \xi_ {1} e _ {1} - \xi_ {2} e _ {2}) = 0.
\]

According to (1) these conditions are written

\[
\left\langle e _ {1}, f (x) \right\rangle = \xi_ {2} \alpha_ {1}, \left\langle e _ {2}, f (x) \right\rangle = - \xi_ {1} \alpha_ {1}.
\]

But we know (loc. cit.) that the image of \(f(E)\) under any linear form on \(E^{*}\) is contained in the ideal \(A\alpha_{1}\) , in other words all the values \(\Phi(x,y)=\langle x,f(y)\rangle\) belong to \(A\alpha_{1}\) ; whence the existence and uniqueness of \(\xi_{1}\) and \(\xi_{2}\) . Thus \(P^{0}\) is a free module of rank \(n-2\) ; there therefore exists, in \(P^{0}\) , by the induction hypothesis, a basis \((e_{3},e_{4},\ldots,e_{n})\) satisfying the conditions of the statement. To show that the basis \((e_{1},\ldots,e_{n})\) of E thus obtained also satisfies these conditions, it suffices to prove that \(\alpha_{1}\) divides \(\alpha_{2}\) ; now this follows from the fact that all the values \(\Phi(x,y)\) are multiples of \(\alpha_{1}\) . It is then clear that \(e_{2r+1},\ldots,e_{n}\) generate \(E^{0}\) . Finally, if \((e_{i}^{\prime})\) is the dual basis of \((e_{i})\) , one has \(f(e_{2j-1})=-\alpha_{j}e_{2j}^{\prime}\) and \(f(e_{2j})=\alpha_{j}e_{2j}^{\prime}\) for \(j=1,\ldots,r\) and \(f(e_{k})=0\) for \(k=2r+1,\ldots,n\) ; the ideals \(A\alpha_{1},A\alpha_{1},A\alpha_{2},A\alpha_{2},\ldots,A\alpha_{r},A\alpha_{r}\) are therefore the invariant factors of \(f(E)\) with respect to \(E^{*}\) , which proves their uniqueness (chap. VII, § 4, no. 2, th. 1).

COROLLARY 1. --- Let A be a commutative field, E a vector space of finite dimension n over A, and Φ an alternating bilinear form on E. Then there exists a basis \((e_i)_{1\leqslant i\leqslant n}\) of E and an even integer \(2r\leqslant n\) such that

\[
\Phi (\sum_ {i = 1} ^ {n} \xi_ {i} e _ {i}, \sum_ {i = 1} ^ {n} \eta_ {i} e _ {i}) = \sum_ {j = 1} ^ {r} (\xi_ {2 j - 1} \eta_ {2 j} - \xi_ {2 j} \eta_ {2 j - 1}).\tag{2}
\]

In particular \(\Phi\) is of even rank \(2r\) .

A basis satisfying (2) is called a symplectic basis for \(\Phi\) .

Remark. --- Note that this corollary is also an immediate consequence of prop. 3 of § 4, no. 2 and of its cor. 1, for with the notations of that proposition, one necessarily has H = \{0\} since Φ is alternating.

COROLLARY 2. --- Let A be a commutative field, E a vector space of finite dimension n over A. For every bivector \(z \in \bigwedge^{2} \mathrm{E}\) there exists a basis \((e_i)_{1 \leqslant i \leqslant n}\) of E such that

\[
z = e _ {1} \wedge e _ {2} + e _ {3} \wedge e _ {4} + \dots + e _ {2 r - 1} \wedge e _ {2 r} \quad (2 r \leqslant n).
\]

It suffices indeed to note that z is canonically identified with an alternating bilinear form on E* (chap. III, § 8, no. 2), and to apply cor. 1 to this form.

Translating cor. 1 into matrix language, we obtain:

COROLLARY 3. --- Let A be a commutative field, R an alternating square matrix over A. The rank of R is an even number 2r, and there exists an invertible matrix P over A such that

\[
{ } ^ { t } P . R . P = \left( \begin{array} { c c c } 0 & I _ { r } & 0 \\ - I _ { r } & 0 & 0 \\ 0 & 0 & 0 \end{array} \right) .
\]

Remark. --- If A is any commutative ring and R is an alternating square matrix of odd order n over A, then det R = 0. This follows from Cor. 3 when A is a field. A direct proof in the case where A is a field of characteristic ≠ 2 is the following: since \(^tR = -R\) , we have det \(R = \det ^t R = (-1)^n\) det R, hence 2 det \(^tR = 0\) . This being so, since the determinant of an alternating matrix ( \(\alpha_{ij}\) ) is a polynomial with integer coefficients in the \(\alpha_{ij}\) such that i \textless{} j, the principle of extension of algebraic identities (chap. IV, § 2, no. 5, Scholium) then shows that our assertion is true for an arbitrary commutative ring A.

\subsection{Pfaffian of a skew-symmetric matrix.}

Let A be a commutative field of characteristic 0, and \(R = (\alpha_{ij})\) an alternating matrix of even order \(2m\) on A. Let E denote the vector space \(\mathbf{A}^{2m}\) , \((e_i)\) ( \(i = 1, \ldots, 2m\) ) its canonical basis, and \(e\) the element \(e_1 \wedge e_2 \wedge \ldots \wedge e_{2m}\) of \(\bigwedge^{2m}\) E. The \(m\) -th exterior power of the bivector \(u = \sum_{i < j} \alpha_{ij} e_i \wedge e_j \in \bigwedge^2 E\) is of the form \(\alpha.e\) , where \(\alpha\) is an element of A that we will compute. The element \(\bigwedge^m u\) is a sum of terms of the form

\[
\alpha_ {h _ {1} k _ {1}} \alpha_ {h _ {2} k _ {2}} \dots \alpha_ {h _ {m} k _ {m}} e _ {h _ {1}} \wedge e _ {k _ {1}} \wedge e _ {h _ {2}} \wedge e _ {k _ {2}} \wedge \dots \wedge e _ {h _ {m}} \wedge e _ {k _ {m}}\tag{3}
\]

with \(h_j < k_j\) for \(j = 1, \ldots, m\) . Such a term is zero if it contains two equal \(e_j\) , that is, if the set \(\{h_1, k_1, \ldots, h_m, k_m\}\) is not exactly \(\{1, 2, \ldots, 2m\}\) . Moreover, if, in (3), one simultaneously interchanges \(e_{h_r}\) and \(e_{h_{r+1}}\) on the one hand, \(e_{k_r}\) and \(e_{k_{r+1}}\) on the other hand, the product does not change; hence it does not change under any permutation performed on the pairs \((h_1, k_1), \ldots, (h_m, k_m)\) . Consider then the sets (and not the sequences) \(S = \{(h_1, k_1), \ldots, (h_m, k_m)\}\) of pairs \((h_j, k_j)\) such that \(1 \leqslant h_j < k_j \leqslant 2m\) for \(j = 1, 2, \ldots, m\) ; let \(\mathscr{F}\) be the set of these pairs. For \(S \in \mathscr{F}\) , set

\[
1 ^ {0}) \varepsilon (S) = 0 \text {   si   } \left\{h _ {1}, k _ {1}, \dots , h _ {m}, k _ {m} \right\} \neq \left\{1, 2, \dots , 2 m \right\};
\]

2°) in the contrary case, \(\varepsilon(S)=1\) or \(\varepsilon(S)=-1\) depending on whether the permutation that maps \(h_{j}\) to 2j-1 and \(k_{j}\) to 2j ( \(j=1,\ldots,m\) ) is even or odd.

The preceding remarks then prove that \(\overset{m}{\wedge} u\) is equal to

\[
m! \sum_ {S \in \mathfrak {S}} \varepsilon (S) (\prod_ {(h, k) \in S} \alpha_ {h k}) e.\tag{4}
\]

We now introduce \(m(2m-1)\) indeterminates \(X_{hk}\) indexed by the pairs \((h,k)\) such that \(1\leqslant h<k\leqslant2m\) and let us call P the polynomial over \(\mathbf{Z}\) with respect to the \(X_{hk}\) defined by

\[
\mathrm{P} ((X _ {h k})) = \sum_ {\mathbf {S} \in \mathfrak {S}} \varepsilon (\mathrm{S}) (\prod_ {(h, k) \in \mathrm{S}} X _ {h k}).\tag{5}
\]

We therefore have

\[
\stackrel {m} {\wedge} u = m! \mathrm{P} ((\alpha_ {h k})). e.\tag{6}
\]

DEFINITION 1. --- Given an alternating matrix \(R = (\alpha_{ij}) (i, j = 1, \ldots, 2m)\) of even order \(2m\) on an arbitrary commutative ring A, one calls the Pfaffian of \(R\) , and denotes it \(\mathrm{Pf}(R)\) , the element \(\mathrm{P}((\alpha_{hk}))\) of A, where \(1 \leqslant h < k \leqslant 2m\) .

Example. --- Suppose:

\[
\alpha_ {1 2} = - \alpha_ {2 1} = \beta_ {1}, \alpha_ {3 4} = - \alpha_ {4 3} = \beta_ {2}, \dots , \alpha_ {2 m - 1, 2 m} = - \alpha_ {2 m, 2 m - 1} = \beta_ {m},
\]

all the others \(\alpha_{ij}\) being zero (cf. th. 1). Then the Pfaffian of \(R = (\alpha_{ij})\) is \(\beta_1\beta_2\ldots\beta_m\) .

PROPOSITION 1. --- Let R be an alternating matrix of even order 2m over a commutative ring A, and let P be a square matrix of order 2m over A. Then

\[
\operatorname{Pf} (^ {t} P. R. P) = (\det P) \operatorname{Pf} (R).\tag{7}
\]

Indeed, suppose first that A is a field of characteristic 0 and set \(R = (\alpha_{ij})\) , \(P = (\beta_{st})\) . Associate with \(R\) the bivector

\[
u = \sum_ {i <   j} \alpha_ {i j} e _ {i} \wedge e _ {j} = \frac {1}{2} \sum_ {1 \leqslant i, j \leqslant 2 m} \alpha_ {i j} e _ {i} \wedge e _ {j}
\]

of \(\wedge\) A \(^{2m}\) , where (e \(_i\) ) denotes the canonical basis of A \(^{2m}\) ; consider \(^t P\) as the matrix, with respect to the basis (e \(_i\) ), of an endomorphism \(f\) of A \(^{2m}\) . Then the bivector ( \(\overset{2}{\wedge} f\) )(u) is associated with the matrix \(^t P.R.P\) since it is equal to \(\frac{1}{2}\sum_{i,j,s,t}\beta_{is}\alpha_{ij}\beta_{jt}e_s\wedge e_t\) . Since the extension \(\wedge f\) of \(f\) to the exterior algebra \(\wedge\) A \(^{2m}\) is an endomorphism of this algebra (chap. III, § 5, no. 9), we have \(\overset{m}{\wedge}((\overset{2}{\wedge}f)(u)) = (\overset{2m}{\wedge}f)(\overset{m}{\wedge}u)\) ; since \(\overset{2m}{\wedge} f\) is the homothety with ratio det \(f\) , it therefore follows from (6) and from def. 1 that we have \(m!\) Pf( \(^t P.R.P\) ) = \(m!\) (det \(P\) )Pf( \(R\) ), whence (7) in the case under consideration. The general case follows by noting that both sides of (7) are polynomials with integer coefficients in the entries of the matrices \(R\) and \(P\) (chap. IV, § 2, no. 5, Scholium).

PROPOSITION 2. --- For every alternating matrix R of even order 2m over a commutative ring A, we have

\[
\det R = (\operatorname{Pf} (R)) ^ {2}.\tag{8}
\]

Indeed, since both sides of (8) are polynomials with integer coefficients in the elements of \(R\) , the principle of extension of algebraic identities (chap. IV, § 2, no. 5, Scholium) shows that it suffices to carry out the proof in the case where A is a field of characteristic 0 and where det \(R \neq 0\) . If \(P\) is an invertible square matrix of order \(2m\) on A, we have det( \(^{t}P.R.P\) ) = ( \(\det P\) ) \(^{2}\) ( \(\det R\) ) and \(\mathrm{Pf}(^{t}P.R.P)\) = ( \(\det P\) ) \(\mathrm{Pf}(R)\) (prop. 1), so that it suffices to prove (8) for \(^{t}P.R.P\) instead of \(R\) . By cor. 1 of th. 1, one can, by a suitable choice of \(P\) , suppose that the matrix \(^{t}P.R.P\) is of the form ( \(\alpha_{ij}\) ) where

\[
\alpha_ {1 2} = - \alpha_ {2 1} = 1, \dots , \alpha_ {2 m - 1, 2 m} = - \alpha_ {2 m, 2 m - 1} = 1,
\]

all the others \(\alpha_{ij}\) being zero (cf. Example). Now the determinant of this matrix equals 1, and its Pfaffian likewise; this completes the proof.

\subsection{Symplectic group.}

Suppose the ring A is commutative. If \(\Phi\) is an alternating bilinear form on E, the automorphisms of the module E leaving \(\Phi\) invariant are called symplectic automorphisms (or symplectic transformations) relative to \(\Phi\) , and they form a group called the symplectic group associated with \(\Phi\) ; it is sometimes denoted \(\mathbf{Sp}(\Phi)\) .

Consider in particular, on the module \(E = A^{2m}\) the alternating bilinear form \(\Phi_{0}\) whose matrix with respect to the canonical basis \((e_{i})\) of E is

\[
R _ {m} = \left( \begin{array}{c c} 0 & I _ {m} \\ - I _ {m} & 0 \end{array} \right).
\]

The symplectic automorphisms and the symplectic group relative to \(\Phi_{0}\) are simply called symplectic automorphisms and symplectic group in \(2m\) variables (over A), respectively; this group is denoted \(\mathbf{Sp}(2m, A)\) or \(\mathbf{Sp}_{2m}(A)\) . Every matrix \(A\) of a symplectic automorphism with respect to the canonical basis \((e)_i\) is called a symplectic matrix. Such a matrix is invertible, and, by formula (48) of § 1, no. 10, satisfies the relation

\[
{ } ^ { t } A . R _ { m } . A = R _ { m } .\tag{9}
\]

Conversely, if a square matrix A of order 2m over A satisfies (9), it is symplectic: it suffices indeed to prove that it is invertible; but (9) implies Pf(R \(_{m}\) ) = Pf( \(^{t}\) A . R \(_{m}\) . A) = (det A)Pf(R \(_{m}\) ) by prop. 1 of no. 2, hence det A = 1. We have at the same time proved the following proposition:

PROPOSITION 3. --- The determinant of a symplectic matrix is equal to 1.

If A is a commutative field and \(\Phi\) a nondegenerate alternating bilinear form on a vector space E of even dimension \(2m\) over A, the symplectic group associated with \(\Phi\) is isomorphic to \(\mathbf{Sp}(2m, \mathrm{A})\) , by cor. 1 of th. 1.

Exercises. --- ¶ 1) Let A be a principal commutative ring, E a free A-module of finite dimension n, Φ an alternating bilinear form on E; the ideals Aαi (1 ≤ i ≤ r) defined in th. 1 of no. 1 are called the invariant factors of Φ.

\begin{enumerate}
\def\labelenumi{\alph{enumi})}
\item
  Let F be a submodule of E, and let \(\Phi_{\mathrm{F}}\) the restriction of \(\Phi\) to F × F. Show that if A \(\beta_{i}\) ( \(1 \leqslant i \leqslant s\) ) are the invariant factors of \(\Phi_{\mathrm{F}}\) (where \(\beta_{i}\) divides \(\beta_{i+1}\) ), we have \(s \leqslant r\) and \(\beta_{i}\) is a multiple of \(\alpha_{i}\) for \(1 \leqslant i \leqslant s\) . (Reduce to the case where \(r = s = n/2\) , and use exerc. 9 b) and 9 c) of chap. VII, § 4).
\item
  Let \(E_{1}\) be a second free A-module of finite dimension, \(\Phi_{1}\) an alternating bilinear form on \(E_{1}, A\gamma_{1}, \ldots, A\gamma_{s}\) its invariant factors ( \(\gamma_{i}\) dividing \(\gamma_{i+1}\) ). In order that \(\Phi_{1}\) be the inverse image of \(\Phi\) by a linear mapping of \(E_{1}\) into E, it is necessary and sufficient that \(s \leqslant r\) and \(\gamma_{i}\) be a multiple of \(\alpha_{i}\) for \(1 \leqslant i \leqslant s\) . (Use a) and prop. 4 of chap. VII, § 4, no. 5).
\item
  Let F, G be two submodules of E, such that \(F^{0}\) (resp. \(G^{0}\) ) be a supplement of F (resp. G) in E. If the restrictions of \(\Phi\) to F and G are equivalent, show that the same holds for the restrictions of \(\Phi\) to \(F^{0}\) and to \(G^{0}\) , and that there exists an automorphism of E leaving \(\Phi\) invariant and mapping F to G.
\item
  Give an example of two submodules F, G of E, of dimension 2, such that F and G admit supplements in E and such that the restrictions \(\Phi_{\mathrm{F}}\) and \(\Phi_{\mathrm{G}}\) are equivalent, but that there exists no automorphism of E leaving \(\Phi\) invariant and mapping F to G (take \(n = 4\) ).
\end{enumerate}

\begin{enumerate}
\def\labelenumi{\arabic{enumi})}
\setcounter{enumi}{1}
\tightlist
\item
  Let \(\Phi\) be an alternating bilinear form on a finite-dimensional vector space E. Show that for every subspace M of E, the difference dim M - dim (M ∩ M⁰) is even. (First consider the case where Φ is nondegenerate).
\end{enumerate}

¶ 3) Let E be a vector space of even dimension n = 2m over a commutative field A; let Φ and Ψ be two alternating bilinear forms on E; suppose Ψ is nondegenerate. Let u and v be the linear maps from E to E* associated on the right with Φ and Ψ respectively; v is an isomorphism from E onto E*; set w = v⁻¹ ∘ u, so that w is an endomorphism of E.

\begin{enumerate}
\def\labelenumi{\alph{enumi})}
\item
  Set \(\mathbf{M}_{\mathbf{0}} = \mathbf{E}\) and by induction \(\mathbf{M}_{k+1} = w(\mathbf{M}_{k})\) for \(k > 0\) . Show that if \(\mathbf{M}_{k}^{\prime}\) is the orthogonal subspace to \(\mathbf{M}_{k}\) for \(\Phi\) , \(\mathbf{M}_{k+1}\) is the orthogonal subspace to \(\mathbf{M}_{k}^{\prime}\) for \(\Psi\) .
\item
  Let \(n_{0}\) be the dimension of \(\mathbf{M}_{0}^{\prime}\) ; we set \(m_{0}=0\) , and for \(k\geqslant1\) , we denote by \(m_{k}\) the dimension of \(\mathbf{M}_{k}\cap\mathbf{M}_{0}^{\prime}\) . Show that for \(k\geqslant1\) the dimension of \(\mathbf{M}_{k}\) is \(n-n_{0}-(m_{1}+\ldots+m_{k-1})\) and that of \(\mathbf{M}_{k}^{\prime}\) is \(n_{0}+(m_{1}+\ldots+m_{k})\) .
\item
  Show that, for all \(k \geqslant 0\) , the dimension of \(\mathbf{M}_k \cap \mathbf{M}_k'\) is \(m_k + m_{k+1} + \ldots + m_{2k}\) , and that of \(\mathbf{M}_k' \cap \mathbf{M}_{k+1}\) is \(m_{k+1} + m_{k+2} + \ldots + m_{2k+1}\) (apply exerc. 2 b) of § 3 to compute the dimensions of \(\mathbf{M}_h \cap \mathbf{M}_{2k-h}'\) and of \(\mathbf{M}_h' \cap \mathbf{M}_{2k+1-h}\) by induction on \(h\) ).
\item
  Deduce from c) that the numbers \(m_{k}\) are even (use exerc. 2).
\item
  Conclude from d) that the number of elementary divisors of the endomorphism w, corresponding to the characteristic root \(\lambda = 0\) , and having a given degree, is even (cf. chap. VII, § 5, exerc. 20).
\end{enumerate}

\begin{enumerate}
\def\labelenumi{\arabic{enumi})}
\setcounter{enumi}{3}
\item
  Let E be a vector space over a commutative field A, admitting a countable basis \((e_n)_{n\geqslant 1}\) , \(\Phi\) a nondegenerate alternating form on E. Show that there exists in E a basis \((a_n)\) such that \(\Phi(a_{2n-1}, a_{2n}) = 1\) for all \(n\geqslant 1\) , and \(\Phi(a_i, a_j) = 0\) for any other pair of indices such that \(i < j\) (reason as in exercise 13 of § 4).
\item
  For any alternating matrix \(X = (x_{ij})\) of even order \(n = 2m\) over a commutative ring, and for every index \(i\) , show that one has \(\operatorname{Pf}(X) = \sum_{j=1}^{n} (-1)^{i+j-1} \operatorname{Pf}(X_{ij}) x_{ij}\) , where \(X_{ij}\) is the matrix of order \(n - 2\) obtained by deleting from \(X\) the rows and columns with indices \(i\) and \(j\) .
\item
  Let M be a square matrix of order m over a commutative ring. Show that if we set
\end{enumerate}

\[
R = \left( \begin{array}{c c} 0 & M \\ - ^ {t} M & 0 \end{array} \right)
\]

we have \(\operatorname{Pf}(R) = \det M\) (first prove it when \(M\) is invertible, using formula (7)).

\begin{enumerate}
\def\labelenumi{\arabic{enumi})}
\setcounter{enumi}{6}
\item
  Let A be a commutative field, \(\mathfrak{A}(A)\) the set of alternating matrices of order \(2m\) on A; let I be a mapping from \(\mathfrak{A}(A)\) into A such that for every matrix \(R\in\mathfrak{A}(A)\) and for every matrix \(P\) of order \(2m\) over A, one has \(\mathrm{I}(^{t}PRP)=(\det P)^{h}\mathrm{I}(R)\) , where \(h\) is a rational integer. Show that \(\mathrm{I}(R)=c(\mathrm{Pf}R)^{h}\) , where \(c\in\mathrm{A}\) (using formula (7) and th. 1 of no. 1).
\item
  Let \(P\) , \(Q\) be two alternating square matrices of even order \(2m\) on a commutative ring A. Let \(\varphi(X) = \mathrm{Pf}(P - XQ)\) ; show that, if \(Q\) is invertible, one has \(\varphi(Q^{-1}P)=0\) . (First consider the case where A is a field; then deduce from exerc. 3 that the minimal polynomial of the matrix \(Q^{-1}P\) divides \(\varphi(X)\) , by passing to an algebraically closed extension of A, and noting that \(\varphi^{2}\) is the characteristic polynomial of \(Q^{-1}P\) up to a scalar factor).
\end{enumerate}

¶ 9) Let E be a vector space of dimension n over a commutative field A, Φ an alternating bilinear form on E, Ψ a sesquilinear Hermitian form (for an involutive automorphism of A) on E, and P and Q the respective matrices of Φ and Ψ with respect to the same basis of E.

\begin{enumerate}
\def\labelenumi{\alph{enumi})}
\item
  We assume \(\Psi\) nondegenerate. Show that the number of elementary divisors of \(Q^{-1}P\) corresponding to the characteristic root 0 and having the same even degree, is an even number (method of exerc. 3).
\item
  We assume that \(\Phi\) be nondegenerate (which implies that \(n\) is even). Show that the number of elementary divisors of \(P^{-1}Q\) corresponding to the characteristic root 0, and having the same odd degree, is an even number (same method).
\end{enumerate}

\begin{enumerate}
\def\labelenumi{\arabic{enumi})}
\setcounter{enumi}{9}
\tightlist
\item
  Let \(\omega_{2m}(\mathbf{F}_q)\) be the order of the symplectic group \(\mathbf{Sp}(2m, \mathbf{F}_q)\) over the finite field \(\mathbf{F}_q\) . Show that if \(h_{2m}\) is the number of pairs of vectors \((x, y)\) of \(\mathbf{F}_q^{2m}\) such that \(\Phi_0(x, y) = 1\) (notations of no. 3), one has \(\omega_{2m}(\mathbf{F}_q) = h_{2m}\omega_{2m-2}(\mathbf{F}_q)\) ; deduce from this
\end{enumerate}

\[
\omega_ {2 m} (\mathbf {F} _ {q}) = (q ^ {2 m} - 1) q ^ {2 m - 1} (q ^ {2 m - 2} - 1) q ^ {2 m - 3} \dots (q ^ {2} - 1) q.
\]

\begin{enumerate}
\def\labelenumi{\arabic{enumi})}
\setcounter{enumi}{10}
\tightlist
\item
  Let A be a commutative field. Show that every transformation \(u\) belonging to the symplectic group \(\mathbf{Sp}(2m, \mathrm{A})\) is a product of transvections belonging to this group (called symplectic transvections; cf.§ 4, exerc. 6). (Reason by induction on \(m\) , by showing that if \(x, y\) are two non-orthogonal vectors of \(\mathrm{E} = \mathrm{A}^{2m}\) , there exists a product \(\nu\) of symplectic transvections such that \(\nu u\) leaves \(x\) and \(y\) invariant). Deduce a new proof of prop. 3 of no. 3.
\end{enumerate}

*12) Let A be a commutative field, E a vector space of even dimension \(n=2m\) over A, \(\Phi\) a nondegenerate alternating bilinear form on E. Show that, for every similarity \(u\) for the form \(\Phi\) ( \(\S 6\) , no. 5), with multiplier \(\alpha\) , we have det \(u=\alpha^{m}\) (using formula (7)).*

¶ 13) We assume that A is a commutative field of characteristic 0, E a vector space of dimension 2m over A, Φ a nondegenerate alternating bilinear form on E. We identify the inverse form \(\widehat{\Phi}\) of Φ with a bivector \(\Gamma\in\bigwedge E\) , so that for any symplectic basis \((e_{i})_{1\leqslant i\leqslant2m}\) of E (for \(\Phi\) ), indexed such that \(\Phi(e_{i},e_{j})=\Phi(e_{m+i},e_{m+j})=0\) , \(\Phi(e_{i},e_{m+j})=\delta_{ij}(1\leqslant i\leqslant m,1\leqslant j\leqslant m)\) , one has

\[
\Gamma = e _ {1} \wedge e _ {m + 1} + e _ {2} \wedge e _ {m + 2} + \dots + e _ {m} \wedge e _ {2 m}.
\]

We say that a nonzero decomposable \(p\) -vector \(z \in \bigwedge E\) is isotropic (resp. totally isotropic) for \(\Phi\) if the vector subspace \(V_z\) corresponding to \(z\) (chap. III, § 7, no. 3) is isotropic (resp. totally isotropic).

\begin{enumerate}
\def\labelenumi{\alph{enumi})}
\item
  If z is a decomposable p-vector \(\neq0\) , 2r the dimension of a supplement of \(V_{z}\cap V_{2}^{0}\) with respect to \(V_{z}\) , show that \(m-p+r\) is the largest of the integers h such that one has \(z\wedge\Gamma^{h}\neq0\) , denoting by \(\Gamma^{h}\) the h-th power of \(\Gamma\) in the exterior algebra \(\wedge E\) (using prop. 2 of § 4, no. 2).
\item
  Show that every bivector \(x \in \bigwedge E\) can be written in the form \(\lambda \Gamma + x_{1}\) , where \(\lambda\) is a scalar and \(x_{1}\) a linear combination of totally isotropic decomposable bivectors. (Reduce to the case where \(x\) is decomposable, and note that if \((e_{i})\) is a symplectic basis, \((e_{1} + e_{2}) \wedge (e_{m+1} - e_{m+2})\) is totally isotropic).
\item
  If \(p \leqslant m\) , show that every \(p\) -vector \(z \in \hat{\Lambda}\) E can be written \(z = x \wedge \Gamma + z_1\) , where \(x\) is a \((p - 2)\) -vector and \(z_1\) a linear combination of \(p\) -vectors that are decomposable and totally isotropic. (Reduce to the case where \(z\) is decomposable, and reason by induction on \(p\) . One can thus reduce to the case where, \((e_i)\) being a symplectic basis, one has \(z = e_1 \wedge e_2 \wedge \ldots \wedge e_{p-1} \wedge e_{m+p-1}\) , and consider the bivectors \((e_{p-1} + e_{p+i}) \wedge (e_{m+p-1} - e_{m+p+i})\) for \(0 \leqslant i \leqslant m - p\) ).
\item
  For \(1 \leqslant p \leqslant m\) , show that every \((m+p)\) -vector \(z \in \wedge\) E can be written \(z = y \wedge \Gamma^p\) , where \(y\) is a \((m-p)\) -vector. (Reduce to the case where \(z\) is decomposable; if \(2r\) is the dimension of a supplement of \(\mathrm{V}_z \cap \mathrm{V}_z^0\) with respect to \(\mathrm{V}_z\) , distinguish two cases, according as \(r = p\) or \(r > p\) ; in the second case, reason by induction on \(r\) , in the same manner as in \(c)\) . Deduce from this that the map \(y \to y \wedge \Gamma^p\) of \(\wedge^{m-p}\) E into \(\wedge^{m+p}\) E is bijective (note that the two spaces have the same dimension). If y is a decomposable \((m-p)\) -vector, show that \(z = y \wedge \Gamma^{p}\) is decomposable and that \(V_{z} = V_{y}^{0}\) .
\end{enumerate}

\begin{enumerate}
\def\labelenumi{\alph{enumi})}
\setcounter{enumi}{4}
\tightlist
\item
  Deduce from d) that, for \(p \leqslant m\) , the subspace \(\hat{E}\) is the direct sum of the homogeneous component of degree p of the two-sided ideal c generated by \(\Gamma\) in \(\hat{E}\) , and of the subspace \(R_{p}\) of p-vectors \(z_{1}\) such that
\end{enumerate}

\(z_{1} \wedge \Gamma^{m-p+1} = 0\) (for \(z \in \bigwedge E\) , apply \(d\) ) to the \((2m - p + 2)\) -vector \(z \wedge \Gamma^{m-p+1}\) ). Using \(c\) ), prove that \(R_{p}\) is generated by the decomposable totally isotropic \(p\) -vectors, and by using \(d\) ), show

that the map \(x \to x \wedge \Gamma\) of \(\wedge E\) into \(\wedge E\) is injective, and that \(R_p\) is of dimension \(\binom{2m}{p} - \binom{2m}{p-2}\) .

\begin{enumerate}
\def\labelenumi{\alph{enumi})}
\setcounter{enumi}{5}
\tightlist
\item
  Show that if \(p \leqslant m\) , a necessary and sufficient condition for a p-vector z to be of the form \(x \wedge \Gamma\) is that, for every decomposable totally isotropic m-vector u, one has \(z \wedge u = 0\) . (Show that, if this condition is satisfied, and if \(z \in R_{p}\) , one has \(z \wedge y = 0\) for all \((2m - p)\) -vector y; for this, put y in the form \(x \wedge \Gamma^{m-p}\) where x is a p-vector, then apply c) to x, and note that if \(x_{1}\) is a decomposable totally isotropic p-vector, \(x_{1} \wedge \Gamma^{m-p}\) is a decomposable \((2m - p)\) -vector which can be written \(u_{1} \wedge \varphi_{1}\) , where \(u_{1}\) is a decomposable totally isotropic \(m\) -vector).
\end{enumerate}

Let a be the annihilator of the ideal c in \(\wedge\) E; a is the direct sum of the subspaces \(R_{m}\) , \(R_{m-1} \wedge \Gamma, \ldots, R_{1} \wedge \Gamma^{m-1}\) and \(K\Gamma^{m}\) . Show that the annihilator of a in \(\wedge\) E is equal to c (cf.§ 2, exerc. 4 b)).

\begin{enumerate}
\def\labelenumi{\alph{enumi})}
\setcounter{enumi}{6}
\tightlist
\item
  For every symplectic automorphism u of E (for \(\Phi\) ), let \(\bar{u}\) be the canonical extension of u to an automorphism of the algebra \(\wedge\) E (chap. III, § 5, no. 9). Show that the only elements of \(\wedge\) E invariant under all the automorphisms \(\bar{u}\) are the linear combinations of 1, \(\Gamma\) , \(\Gamma^{2}\) , \(\ldots\) , \(\Gamma^{m}\) with coefficients in K. (If ( \(e_{i}\) ) is a symplectic basis, write that an element of \(\wedge\) E is invariant under the symplectic transvections (exerc. 11) corresponding to the hyperplanes orthogonal to the \(e_{i}\) ; then consider the symplectic transformations \(u_{ij}\) defined by \(u_{ij}(e_{i}) = e_{j}\) , \(u_{ij}(e_{j}) = e_{i}\) , \(u_{ij}(e_{m+i}) = e_{m+j}\) , \(u_{ij}(e_{m+j}) = e_{m+i}\) , \(u_{ij}(e_{k}) = e_{k}\) for every other index k).
\end{enumerate}

\begin{enumerate}
\def\labelenumi{\arabic{enumi})}
\setcounter{enumi}{13}
\item
  Let A be a commutative field of characteristic 2, E the vector space \(\mathrm{A}^{2m}\) ; we identify the symplectic group \(\mathbf{Sp}(2m,\mathrm{A})\) with the group of symplectic matrices \(U\) , thus satisfying the relation \(^t U.R.U = R\) , where \(R = \left( \begin{array}{cc}0 & I_m\\ I_m & 0 \end{array} \right)\) . We set \(D = \left( \begin{array}{cc}0 & 0\\ I_m & 0 \end{array} \right)\) ; show that every symplectic matrix \(U\) such that \(I + U\) is invertible can be written in a unique way in the form \((^t D + S)^{-1}(D + S)\) , where \(S\) is a symmetric matrix such that \(D + S\) is invertible, and conversely (cf.§ 4, exerc. 11).
\item
  Let A be a commutative field, E a vector space of dimension 2m over A, Φ a nondegenerate alternating bilinear form on E. a) Extend to endomorphisms u of E such that \(u^{*}u = uu^{*}\) ( \(u^{*}\) being the adjoint of u with respect to Φ) the results of § 4, exerc. 12 a) and b).
\end{enumerate}

\begin{enumerate}
\def\labelenumi{\alph{enumi})}
\setcounter{enumi}{1}
\item
  We assume that \(u^{*} = u\) . With the notations of exerc. 12 of § 4, show that if M is a minimal element of the set \(\mathfrak{M}\) of non-isotropic subspaces contained in \(G(p, p)\) and stable under \(u\) , then either M is an indecomposable submodule of \(E_{u}\) , or it is a direct sum of two such isomorphic submodules (reason as in exerc. 12 c) of § 4). Show by examples (with \(p(X) = X - 1\) ) that the two cases can occur.
\item
  We assume \(u^* = -u\) , with A not of characteristic 2. With the same notations, let \(p^h\) be the minimal polynomial of the restriction of \(u\) to M, and let \(d\) be the degree of \(p\) . Show that if \(d(h - 1)\) is odd, M is an indecomposable submodule of \(E_u\) ; if on the contrary \(d(h - 1)\) is even, M is either indecomposable or a direct sum of two isomorphic indecomposable submodules.
\item
  We assume that \(u^{*}u = 1\) (in other words that \(u \in \mathbf{Sp}(\Phi)\) ). With the notations of c), if \(p(X)\) does not divide \(X^{d(h-1)} - 1\) , or if \(p(X) = X - 1\) and \((-1)^{h} \neq -1\) , M is an indecomposable submodule of \(E_{u}\) ; otherwise, M is either indecomposable or a direct sum of two isomorphic indecomposable submodules.
\end{enumerate}

\section{Special properties of Hermitian forms}

\subsection{Orthogonal bases.}

DEFINITION 1. --- Let Φ be a Hermitian form on E. A basis ( \(e_{i}\) ) of E is said to be orthogonal for Φ if any two elements of this basis are orthogonal for Φ.

If moreover \(\Phi(e_i, e_i) = 1\) for all \(i\) , the basis \((e_i)\) is said to be orthonormal.

Let \((e_i)\) be an orthogonal basis; if we set \(\Phi(e_i, e_i) = \alpha_i\) , one has

\[
\Phi (\sum_ {i} \xi_ {i} e _ {i}, \sum_ {i} \eta_ {i} e _ {i}) = \sum_ {i} \xi_ {i} \alpha_ {i} \overline {{\eta}} _ {i}.\tag{1}
\]

Lemma 1. --- Suppose that A is a field and \(\Phi\) a Hermitian form \(\neq 0\) on E. If all vectors of E are isotropic, then A is a commutative field of characteristic 2, the antiautomorphism J is the identity, and \(\Phi\) is alternating.

Indeed, if one expands \(\Phi(x+y, x+y)=0\) , one obtains, taking the hypotheses into account, \(\Phi(x,x)=\Phi(y,y)=0\) , the relation \(\Phi(x,y)=-\overline{\Phi(x,y)}\) for all \(x,y\) in E. Since \(\Phi\) is \(\neq 0\) , there exist \(x,y\) in E such that \(\Phi(x,y)=1\) . Writing that \(\Phi(\lambda x,y)=-\overline{\Phi(\lambda x,y)}\) , one obtains \(\bar{\lambda}=-\lambda\) for all \(\lambda\in A\) . Taking first \(\lambda=1\) , one sees that A is of characteristic 2; the relation \(\bar{\lambda}=-\lambda\) then shows that J is the identity, hence A is commutative and \(\Phi\) is alternating.

THEOREM 1. --- Suppose that A is a field and E a vector space of finite dimension n over A. Then, for every Hermitian form Φ on E, E admits an orthogonal basis, except if the following conditions are simultaneously satisfied:

\begin{enumerate}
\def\labelenumi{(\Alph{enumi})}
\setcounter{enumi}{2}
\tightlist
\item
  A is commutative of characteristic 2, the antiautomorphism J is the identity, Φ is alternating and nonzero.
\end{enumerate}

Reason by induction on \(n\) , the result being evident for \(n=0\) . We may suppose \(\Phi\neq0\) . If (C) is not satisfied, Lemma 1 shows that there exists an element \(x\in\mathbf{E}\) such that \(\Phi(x,x)\neq0\) . Let H be the subspace of E orthogonal to \(x\) ; it has dimension \(\geqslant n-1\) , and, since \(x \notin H\) , \(H\) is exactly of dimension \(n-1\) . If the restriction \(\Psi\) of \(\Phi\) to \(H\) does not satisfy (C), then by the induction hypothesis there exists a basis \((e_2, \ldots, e_n)\) of \(H\) which is orthogonal for \(\Psi\) ; by setting \(e_1 = x\) one obtains an orthogonal basis \((e_1, e_2, \ldots, e_n)\) of \(E\) . It remains to examine the case where \(A\) is a commutative field of characteristic 2, where \(J\) is the identity, and where \(\Psi\) is alternating and \(\neq 0\) . There then exists \(y\) , \(z\) in \(H\) such that \(\Psi(y, z) \neq 0\) ; set \(e_1 = x + y\) ; for \(x + \lambda z (\lambda \in A)\) to be orthogonal to \(e_1\) , it is necessary and sufficient that we have \(0 = \Phi(x + y, x + \lambda z) = \Phi(x, x) + \lambda \Psi(y, z)\) , a condition which determines \(\lambda\) in a unique way; the scalar \(\lambda\) having been chosen thus, we have \(\Phi(x + \lambda z, x + \lambda z) = \Phi(x, x) \neq 0\) , hence the restriction \(\Psi'\) of \(\Phi\) to the subspace \(H'\) of \(E\) orthogonal to \(e_1\) is not alternating; one may consequently apply the induction hypothesis to \(H'\) , which proves the theorem.

When (C) is satisfied, there evidently exists no orthogonal basis for \(\Phi\) .

COROLLARY 1. --- With the notations of th. 1, one supposes moreover that (C) is not satisfied, that \(\Phi\) is nondegenerate and that, for every \(x \in E\) , there exists \(\rho \in A\) such that \(\Phi(x, x) = \rho \bar{\rho}\) . Then E admits an orthonormal basis for \(\Phi\) .

Indeed, let \((e_i)\) ( \(i = 1, \ldots, n\) ) be an orthogonal basis of E. Put \(\Phi(e_i, e_i) = \alpha_i\) . We have \(\alpha_i \neq 0\) for \(i = 1, \ldots, n\) since \(\Phi\) is nondegenerate. By hypothesis, there exist elements \(\beta_i\) of A such that \(\alpha_i = \beta_i \overline{\beta_i}\) for \(i = 1, \ldots, n\) ; we have \(\beta_i \neq 0\) . By setting \(f_i = \beta_i^{-1} e_i\) , one has \(\Phi(f_i, f_i) = \beta_i^{-1} \alpha_i \overline{\beta_i^{-1}} = 1\) for all \(i\) , and \(\Phi(f_i, f_j) = 0\) for \(i \neq j\) . Hence \((f_i)\) is an orthonormal basis.

Remark. --- The last hypothesis of the corollary is satisfied when J is the identity and every element of A is the square of an element of A (for example, when A is algebraically closed).

COROLLARY 2. --- Let A be a field and R a Hermitian matrix of order n and rank r over A. Then, unless the following condition holds:

(C') A is commutative of characteristic 2, J is the identity, R is alternating and nonzero,

there exists an invertible matrix P of order n over A such that

\[
{ } ^ { t } P . R . \overline { { P } } = \left( \begin{array} { c c c c c } \alpha _ { 1 } & 0 \dots 0 \dots 0 \\ 0 & \alpha _ { 2 } \dots 0 \dots 0 \\ \cdot & \cdot & \cdot & \cdot & \cdot \\ 0 & 0 \dots \alpha _ { r } \dots 0 \\ 0 & 0 \dots 0 \dots 0 \\ \cdot & \cdot & \cdot & \cdot & \cdot \\ 0 & 0 \dots 0 \dots 0 \end{array} \right)
\]

\(where \overline{\alpha}_{i} = \alpha_{i} \neq 0 \text{ for } i = 1, \ldots, r.\)

PROPOSITION 1. --- Suppose that A is a commutative field. Let \(\Phi\) be a Hermitian form on E, and \((x_{n})\) ( \(n=1,2,\ldots\) ) a sequence (finite or infinite) of linearly independent vectors of E such that, for every \(n\) , the subspace \(\mathrm{E}_{n}=\mathrm{A}x_{1}+\cdots+\mathrm{A}x_{n}\) is non-isotropic. Let \(\mathrm{D}_{jn}\) ( \(j\leqslant n\) ) be the cofactor of \(\Phi(x_{j},x_{n})\) in the matrix \((\Phi(x_{s},x_{t}))_{(s,t=1,\ldots,n)}\) . We then have \(\mathrm{D}_{nn}\neq0\) for all \(n\) . Set

\[
e _ {n} = \sum_ {j = 1} ^ {n} \mathrm{D} _ {n n} ^ {- 1} \mathrm{D} _ {j n} x _ {j}.\tag{2}
\]

Then, for every \(n\) , \((e_{1},\ldots,e_{n})\) is an orthogonal basis of \(\mathrm{E}_{n}\) and we have

\[
\Phi (e _ {n}, e _ {n}) = \mathrm{D} _ {n n} ^ {- 1} \mathrm{D} _ {n + 1, n + 1}.\tag{3}
\]

Indeed, since the restriction of \(\Phi\) to \(\mathbf{E}_{n-1}\) is nondegenerate, we have \(\mathrm{D}_{nn} \neq 0\) ( \(\S 2\) , prop. 3); note that we have \(\mathrm{D}_{11} = 1\) since the determinant of the empty matrix is equal to 1. Formulas (2) imply first that we have \(e_n \equiv x_n\) (mod. \(\mathbf{E}_{n-1}\) ) for all \(n\) , hence that the \(e_n\) are linearly independent, and that \((e_1, \ldots, e_n)\) is a basis of \(\mathbf{E}_n\) . For every \(j < n\) , one has

\[
\Phi (e _ {n}, x _ {j}) = \mathrm{D} _ {n n} ^ {- 1} \sum_ {k = 1} ^ {n} \mathrm{D} _ {k n} \Phi (x _ {k}, x _ {j}) = 0
\]

(chap. III, § 6, no. 1, formula (12)); hence \(e_n\) is orthogonal to \(\mathrm{E}_{n-1}\) , and in particular to \(e_j\) for \(j < n\) . On the other hand, we have

\[
\begin{array}{r l} \Phi (e _ {n}, e _ {n}) & = \Phi (e _ {n}, \sum_ {j = 1} ^ {n} \mathrm{D} _ {n n} ^ {- 1} \mathrm{D} _ {j n} x _ {j}) = \Phi (e _ {n}, x _ {n}) = \Phi (\sum_ {j = 1} ^ {n} \mathrm{D} _ {n n} ^ {- 1} \mathrm{D} _ {j n} x _ {j}, x _ {n}) \\ & = \mathrm{D} _ {n n} ^ {- 1} \sum_ {j = 1} ^ {n} \mathrm{D} _ {j n} \Phi (x _ {j}, x _ {n}) = \mathrm{D} _ {n n} ^ {- 1} \mathrm{D} _ {n + 1, n + 1} \end{array}
\]

(chap. III, § 6, no. 1, formula (10)). This proves our assertions.

With the notations of prop. 1, one says that the sequence \((e_n)\) is obtained from the sequence \((x_n)\) by the Gram-Schmidt orthogonalization process.

PROPOSITION 2. --- Let \(\Phi\) be a Hermitian form on E, and \((e_i)\) ( \(i = 1, \ldots, n\) ) an orthogonal basis (resp. orthonormal basis) of E for \(\Phi\) . Then, for every \(p \geqslant 0\) , the basis of \(\bigotimes^{p} \mathbf{E}\) formed by the \(e_{i_1} \otimes \cdots \otimes e_{i_p}\) and the basis \((e_h)\) of \(\bigwedge^{p} \mathbf{E}\) (where H ranges over the set of subsets with \(p\) elements of {[}1, n{]}; cf. chap. III, § 5, no. 6) are orthogonal (resp. orthonormal) for the extensions of \(\Phi\) to \(\bigotimes^{p} \mathbf{E}\) and \(\bigwedge^{p} \mathbf{E}\) respectively (§ 1, no. 9). If, moreover, the maps associated with \(\Phi\) are bijective, the basis \((e_i')\) of \(\mathbf{E}^*\) dual of \((e_i)\) is orthogonal (resp. orthonormal) for the inverse form \(\hat{\Phi}\) of \(\Phi\) (§ 1, no. 7).

The assertions concerning \(\widehat{\otimes}\) E and \(\wedge\) E follow immediately from formulas (35) and (37) of § 1, no. 9. The one relating to the inverse form follows from the fact that the matrix of \(\widehat{\Phi}\) with respect to \((e_i')\) is the inverse of the matrix of \(\Phi\) with respect to \((e_i)\) ( \(§ 1\) , no. 10).

\subsection{Unitary group and orthogonal group.}

Let \(\Phi\) be a Hermitian form on E; the automorphisms of the A-module E that leave \(\Phi\) invariant are called unitary automorphisms (or unitary transformations) relative to \(\Phi\) , and their group is called the unitary group associated with \(\Phi\) ; it is denoted \(\mathbf{U}(\Phi)\) . Given a quadratic form \(Q \neq 0\) on E, the automorphisms of the A-module E which leave Q invariant are called orthogonal automorphisms (or orthogonal transformations) relative to Q; their group is called the orthogonal group associated with Q; it is denoted by \(\mathbf{O}(Q)\) .

Every orthogonal transformation for a quadratic form Q is unitary for the bilinear form associated with Q. The converse is true when the scalar 2 is not equal to 0 or a zero divisor in A (§ 3, no. 4, (13)), for example if A is a field of characteristic ≠ 2.

Consider, in particular, on the module \(E = A^{n}\) the Hermitian form \(\Phi_{0}\) whose matrix with respect to the canonical basis \((e_{i})\) of E is the identity matrix \(I_{n}\) . The unitary automorphisms associated with \(\Phi_{0}\) are simply called unitary automorphisms (or transformations) in n variables; their group is called the unitary group in n variables and is sometimes denoted \(\mathbf{U}(n, \mathbf{A})\) or \(\mathbf{U}_{n}(\mathbf{A})\) . The matrix U of a unitary automorphism with respect to \((e_{i})\) is called a unitary matrix. Such a matrix is invertible, and satisfies, according to formula (48) of § 1, no. 10, the relation

\[
{ } ^ { t } U . \overline { { U } } = I _ { n } ;\tag{4}
\]

Conversely, if A is a commutative ring or a field, a matrix U satisfying (4) is invertible, and is then unitary.

When J is the identity and 2 is neither equal to 0 nor a divisor of 0 in A, one uses the terms orthogonal group in n variables, orthogonal automorphism (or orthogonal transformation) in n variables, and orthogonal matrix instead of the preceding terms, and one writes \(\mathbf{O}(n,\mathrm{A})\) (resp. \(\mathbf{O}_n(\mathrm{A})\) ) instead of \(\mathbf{U}(n,\mathrm{A})\) (resp. \(\mathbf{U}_n(\mathrm{A})\) ). Relation (4) is then written

\[
{ } ^ { t } U . U = I _ { n }\tag{5}
\]

and, since A is commutative, it is a necessary and sufficient condition for U to be an orthogonal matrix.

PROPOSITION 3. --- Suppose that A is a commutative field and that E is of finite dimension \textgreater{} 0. Let Φ be a nondegenerate Hermitian form on E. The map u → det u is a homomorphism from the unitary group U(Φ) associated with Φ onto the multiplicative subgroup H of A consisting of the elements ρ such that ρρ̄ = 1 (a subgroup reduced to \{1, -1\} when J is the identity).

Indeed, let \(u\) be an element of \(\mathbf{U}(\Phi)\) , \(U\) its matrix with respect to a basis of E, and \(R\) the matrix of \(\Phi\) with respect to this basis. The relation \(R = {}^t U.R.\overline{U}\) ( \(\S 1\) , no. 10, formula (48)) shows that one has \((\det U)(\det \overline{U}) = 1\) since \(R\) is invertible; whence \((\det u)(\overline{\det u}) = 1\) . The homomorphism \(u \to \det u\) maps \(\mathbf{U}(\Phi)\) onto H. Indeed, when A is of characteristic 2 and J is the identity, H is reduced to the element 1. Otherwise there exists an orthogonal basis \((e_i)(i = 1,\ldots,n)\) of E (th. 1); for every \(\rho \in A\) such that \(\rho \bar{\rho} = 1\) , let \(u\) be the automorphism of E defined by \(u(e_1) = \rho e_1\) and \(u(e_i) = e_i\) for \(i = 2,\ldots,n\) ; then \(u\) is unitary and det \(u = \rho\) , whence the proposition.

Under the conditions of prop. 3, the kernel of the homomorphism \(u \to \det u\) is a normal subgroup of \(\mathbf{U}(\Phi)\) , which is called the special unitary group associated with \(\Phi\) ; it is sometimes denoted \(\mathbf{SU}(\Phi)\) .

When J is the identity and A is not of characteristic 2, this group is also called the special orthogonal group associated with \(\Phi\) (or with the quadratic form \(Q(x) = \Phi(x, x)\) ) and is sometimes denoted SO(Q).

If E = A \(^{n}\) and if Φ is the form whose matrix with respect to the canonical basis of E is the identity matrix, one uses the notations SU(n, A) or SU \(_{n}\) (A) and SO(n, A) or SO \(_{n}\) (A).

\subsection{Orthogonal projectors and involutions.}

Throughout this no., one assumes that the scalar 2 is \(invertible\) in A (for example that A is a field of characteristic \(\neq 2\) ), and that \(\Phi\) is a Hermitian form \(nondegenerate\) on E. One denotes by \(\frac{1}{2}\) the inverse of 2.

Lemma 2. --- For an endomorphism u of E to be such that \(u^{2}=1\) , it is necessary and sufficient that \(\frac{1}{2}(1-u)\) be a projector in E; then u is the difference of the two projectors \(\frac{1}{2}(1+u)\) and \(\frac{1}{2}(1-u)\) .

Indeed, in the ring \(\mathfrak{L}(\mathrm{E})\) , the relation \(\left(\frac{1}{2}(1 - u)\right)^2 = \frac{1}{2}(1 - u)\) is equivalent to \(u^2 = 1\) . The rest is trivial.

An endomorphism \(u\) of E such that \(u^2 = 1\) (which is then necessarily an automorphism of E equal to its inverse) is called an involution. Set \(v = \frac{1}{2}(1 - u)\) , \(U^- = v(E)\) , \(U^+ = v^{-1}(0)\) ( \(= w(E)\) by setting \(w = \frac{1}{2}(1 + u)\) ); one knows that E is the direct sum of \(U^+\) and of \(U^-\) (chap. VIII, § 1, no. 1), and one has \(u(x) = x\) in \(U^+\) , \(u(x) = -x\) in \(U^-\) . When A is a field and E is of finite dimension, it follows, since A is of characteristic \(\neq 2\) , that the only eigenvectors \(\neq 0\) of \(u\) are the elements ≠ 0 in U⁺ or in U⁻; they correspond respectively to the eigenvalues +1 and −1.

PROPOSITION 4. --- Let \(u \in \mathbf{GL}(\mathrm{E})\) be an involution. The following properties are equivalent:

\begin{enumerate}
\def\labelenumi{\alph{enumi})}
\item
  u belongs to the unitary group associated with \(\Phi\) ;
\item
  the submodules \(U^{+} = \frac{1}{2}(1 + u)(E)\) and \(U^{-} = \frac{1}{2}(1 - u)(E)\) are orthogonal (and consequently non-isotropic).
\end{enumerate}

Moreover, if A is a field and E is finite-dimensional, properties a) and b) are equivalent to:

\begin{enumerate}
\def\labelenumi{\alph{enumi})}
\setcounter{enumi}{2}
\tightlist
\item
  \(u = u^{*}\) .
\end{enumerate}

Indeed, for \(x \in \mathrm{U}^{+}\) and \(y \in \mathrm{U}^{-}\) , the relation \(\Phi(u(x), u(y)) = \Phi(x, y)\) gives \(2\Phi(x, y) = 0\) , hence \(a)\) implies \(b)\) . Conversely, we obviously have \(\Phi(u(x), u(y)) = \Phi(x, y)\) when \(x\) and \(y\) are both in \(\mathrm{U}^{+}\) or both in \(\mathrm{U}^{-}\) , and, given \(b)\) , this relation is still true when one of them is in \(\mathrm{U}^{+}\) and the other in \(\mathrm{U}^{-}\) ; since E is the direct sum of \(\mathrm{U}^{+}\) and \(\mathrm{U}^{-}\) , we see that \(b)\) implies \(a)\) . Finally, when E is a finite-dimensional vector space, the adjoint \(u^{*}\) is defined since \(\Phi\) is nondegenerate; the relation \(a)\) is equivalent to \(uu^{*} = 1\) ( \(§ 1, n^{\circ} 8\) , cor. of prop. 8); since \(u^{2} = 1\) by hypothesis, \(a)\) and \(c)\) are equivalent.

COROLLARY 1. --- Suppose that A is a field and that E is finite-dimensional. The map \(u \to \frac{1}{2}(1 + u)(E)\) is a bijection from the set of involutions u belonging to the unitary group associated with \(\Phi\) onto the set of non-isotropic subspaces of E; the subspace U \(^+\) corresponding to u is the set of elements of E invariant under u.

By prop. 4, it suffices to show that every non-isotropic subspace M of E is the set of vectors invariant under an involution \(u \in \mathbf{U}(\Phi)\) , and that this involution is unique. Now E is the direct sum of M and of \(\mathbf{M}^0\) ( \(\S 4\) , no. 1, cor. of prop. 1), and one necessarily has \(u(x) = x\) for \(x \in \mathbf{M}\) and \(u(x) = -x\) for \(x \in \mathbf{M}^0\) by virtue of prop. 4; these relations determine \(u\) uniquely, and the endomorphism \(u\) thus determined evidently answers the question (prop. 4).

One says that the involution \(u\) thus determined is the symmetry with respect to the non-isotropic subspace M.

COROLLARY 2. --- Suppose that A is a field and that E is of finite dimension. For a projector v in E to be such that v(E) and \(v^{-1}(0)\) are orthogonal (and consequently non-isotropic), it is necessary and sufficient that v = v*.

It suffices to apply prop. 4 to the involution \(u = 1 - 2v\) .

A projector satisfying the condition of corollary 2 is called an orthogonal projector for \(\Phi\) .

\subsection{Symmetries in the orthogonal group.}

Unless expressly stated otherwise, one assumes, in this no., that A is a commutative field of characteristic ≠ 2, and that Φ is the symmetric bilinear form associated with a nondegenerate quadratic form Q on E. Recall that one has Φ(x, x) = 2Q(x) for x ∈ E (§ 3, no. 4).

Let H be a non-isotropic hyperplane in E, and u the symmetry with respect to H (no. 3). Let \(a \neq 0\) be a vector orthogonal to H; one has by hypothesis \(u(a) = -a\) . Every vector \(x \in E\) can be written in one and only one way in the form \(x = \lambda a + y\) with \(\lambda \in A\) and \(y \in H\) ; since \(a\) and \(y\) are orthogonal, one has \(\Phi(x, a) = \lambda \Phi(a, a)\) , whence, since \(a\) is non-isotropic ( \(\S 4\) , no. 1, cor. of prop. 1), \(\lambda = \Phi(x, a) \Phi(a, a)^{-1}\) . This being so, one has

\[
u (x) = \lambda u (a) + u (y) = - \lambda a + y = x - 2 \lambda a,
\]

whence

\[
u (x) = x - 2 \Phi (x, a) \Phi (a, a) ^ {- 1}. a = x - \Phi (x, a) Q (a) ^ {- 1}. a.\tag{6}
\]

One notes that the last member of (6) retains a meaning when A is a field of characteristic 2, and a is a non-singular vector of E; one verifies immediately that one still has then Q(u(x)) = Q(x) for all x ∈ E, in other words u ∈ O(Q). One also says that the involution u thus defined is the symmetry with respect to the hyperplane orthogonal to a (cf. exerc. 28).

PROPOSITION 5. --- Suppose the vector space E is of finite dimension n. The orthogonal group \(\mathbf{O}(\mathrm{Q})\) associated with Q is then generated by the symmetries with respect to the non-isotropic hyperplanes of E.

The proposition being evident for \(n=0\) , we shall reason by induction on \(n\) . Let \(u\) be an orthogonal transformation of E, and let \(x\) be a non-isotropic vector of E (Lemma 1); we distinguish three cases:

\begin{enumerate}
\def\labelenumi{\alph{enumi})}
\item
  Suppose first that \(u(x) = x\) . Then the hyperplane H orthogonal to \(x\) is non-isotropic, and we have \(u(\mathrm{H}) = \mathrm{H}\) . The restriction \(u'\) of \(u\) to H therefore belongs to the orthogonal group \(\mathbf{O}(Q')\) associated with the restriction \(Q'\) of Q to H. The induction hypothesis implies, since \(Q'\) is nondegenerate, that one has \(u' = v_1' \ldots v_m'\) , where \(v_i'\) is a symmetry with respect to a hyperplane \(L_i\) of H. The endomorphism \(v_i\) of E which extends \(v_i'\) and is such that \(v_i(x) = x\) is then the symmetry with respect to the hyperplane \(Ax + L_i\) of E. We obviously have \(u = v_1v_2 \ldots v_m\) .
\item
  Suppose secondly that \(u(x) = -x\) . If we denote by \(s\) the symmetry with respect to the hyperplane H orthogonal to \(x\) , and if we set \(\nu = su\) , one has \(\nu(x) = x\) , and we are reduced to the case \(a\) ).
\item
  Let us finally pass to the general case, and set \(y = u(x)\) so that \(Q(y) = Q(x)\) . Under these conditions, the vectors \(x - y\) and \(x + y\) cannot both be isotropic, for from the relations \(Q(x - y) = 0\) and \(Q(x + y) = 0\) , one would deduce, by adding memberwise, \(2(Q(x) + Q(y)) = 0\) (§ 3, no. 4, def. 2), whence \(4Q(x) = 0\) , contrary to the hypothesis. Suppose, for example, that \(a = x - y\) is not isotropic; we then have
\end{enumerate}

\[
\Phi (y, a) = \mathrm{Q} (y + a) - \mathrm{Q} (y) - \mathrm{Q} (a) = \mathrm{Q} (x) - \mathrm{Q} (y) - \mathrm{Q} (a) = - \mathrm{Q} (a);
\]

consequently, if one denotes by \(s\) the symmetry with respect to the hyperplane orthogonal to \(a\) , formula (6) proves that \(s(y) = y + a = x\) ; by setting \(\nu = su\) , one has \(\nu(x) = x\) , and we are reduced to the case \(a\) . If \(a = x - y\) is isotropic and \(b = x + y\) not isotropic; one sees in the same way that we are reduced to the case \(b\) .

\subsection{The group of similarities.}

Let \(\Phi\) be a Hermitian form on E. An automorphism \(u\) of the A-module E is called a similarity (relative to \(\Phi\) ) if there exists an invertible element \(\alpha\) of A such that one has

\[
\Phi (u (x), u (y)) = \alpha \Phi (x, y)\tag{7}
\]

for all x, y in E. The similarities form a group Γ.

When \(\Phi\) takes values that are regular elements of A (for example when A is a field and \(\Phi \neq 0\) ), the element \(\alpha\) of A satisfying (7) is uniquely determined by \(u\) ; it is called the multiplier of the similarity \(u\) . Changing \(x\) to \(\lambda x\) in (7), we then see that \(\alpha\) belongs to the center of A; exchanging \(x\) and \(y\) in (7), we see moreover that \(\overline{\alpha} = \alpha\) . If, for \(u \in \Gamma\) , we denote by \(\alpha(u)\) the multiplier of \(u\) , the map \(u \to \alpha(u)\) is a homomorphism of \(\Gamma\) into the multiplicative group of invertible elements of the center of A. The kernel of this homomorphism is the unitary group associated with \(\Phi\) , which is therefore a normal subgroup of \(\Gamma\) . Let \(\beta\) be an invertible element of the center of A, \(v\) the homothety of ratio \(\beta\) , and \(w\) a unitary automorphism of E; then \(vw = wv\) is a similarity of E, and its multiplier is \(\beta\overline{\beta}\) . Conversely, let \(u\) be a similarity whose multiplier is of the form \(\beta\overline{\beta}\) ( \(\beta\) denoting an invertible element of the center of A); then \(uv^{-1}\) is a unitary automorphism \(w\) , hence \(u\) is of the form \(vw\) .

Suppose now that A is a field, E is a finite-dimensional vector space, and that \(\Phi\) is nondegenerate. For every similarity \(u\) with multiplier \(\alpha\) we have

\[
\Phi (x, \alpha y) = \alpha \Phi (x, y) = \Phi (u (x), u (y)) = \Phi (x, u ^ {*} (u (y))),
\]

therefore \(u^*u\) is the homothety with ratio \(\alpha\) . If A is commutative, and if \(n\) denotes the dimension of E, it follows from this and from formula (50) of § 1, no. 10 that one has

\[
(\det u) \overline {{(\det u)}} = \alpha^ {n}.\tag{8}
\]

Let us then distinguish two cases:

1°) The integer n is odd, i.e., \(n = 2q + 1\) . Then, setting \(\rho = \alpha^{-q}(\det u)\) , one has \(\alpha = (\det u) \overline{(\det u)} \alpha^{-2q} = \rho \bar{\rho}\) . Thus u is the product of the homothety with ratio \(\rho\) and a unitary automorphism.

2°) The integer n is even, say n = 2q. Then, setting \(\rho = \alpha^{-q}(\det u)\) , one has \(\rho\bar{\rho} = 1\) . In particular, when J is the identity, one has \((\det u)^{2} = \alpha(u)^{2q}\) ; the similarities u such that det \(u = \alpha(u)^{q}\) (resp. det \(u = -\alpha(u)^{q}\) ) are called direct (resp. inverse); the direct similarities form a normal subgroup of index 2 of \(\Gamma\) ;

the homotheties with ratio ≠ 0 are direct similarities; the same is true of orthogonal transformations with determinant 1 (no. 2); orthogonal transformations with determinant -1 are inverse similarities.

The preceding definitions and results remain valid for ε-hermitian forms (§ 3, no. 1), and in particular for alternating forms.

Let A be a commutative field and Q a quadratic form \(\neq 0\) on E. One calls a similarity (relative to Q) any automorphism \(u\) of E such that there exists a nonzero element \(\alpha\) of A (called the multiplier of \(u\) ) for which \(Q(u(x)) = \alpha Q(x)\) for every \(x \in E\) . It is clear that \(u\) is then a similarity with multiplier \(\alpha\) relative to the bilinear form associated with Q; the converse is true when the characteristic of A is \(\neq 2\) .
